\documentclass[addpoints, 12pt]{exam}

\usepackage{amsmath}
\usepackage{amssymb}
\usepackage{xifthen} % provides \isempty test
\usepackage{graphicx}
\usepackage[hidelinks]{hyperref}
\usepackage{tabularx}
\usepackage{booktabs}
\usepackage{multirow}
\usepackage{enumitem}
\usepackage{tikz}
\usetikzlibrary{calc}
\usepackage{siunitx}
\usepackage{soul}

\sisetup{
  group-separator = {\,},
  group-minimum-digits = 4
}

% Tightlist environment when using "-" for bullet points.
\setlength{\emergencystretch}{3em} % prevent overfull lines
\providecommand{\tightlist}{%
  \setlength{\itemsep}{0pt}\setlength{\parskip}{0pt}}

% Highlighting code chuncks

\usepackage{setspace}
\onehalfspacing

% CSL referencing

\pagestyle{headandfoot}
\runningheadrule
\firstpageheader{\large \textbf{}}
{\large \textbf{\Large TEST X -- MODULE CODE}}
{\large \textbf{}}
%% Other pages
\runningheader{\textbf{MODULE CODE}}
{\textbf{}}
{\textbf{Page \thepage~of \numpages~}}

\firstpagefooter{\large \textbf{}}
{\large \textbf{}}
{\large \textbf{}}
%% Other pages
\runningfooter{\textbf{}}
{\textbf{}}
{\textbf{}}

\def\tick{
  \begin{tikzpicture}[scale=0.3]
    \draw[fill=black] (0,.35) -- (.25,0) -- (1,.7) -- (.25, .175) -- cycle;
    \node[above, white] at (0.35, 0.1) {\tiny 1};
    \node[below, white] at (0.7, 0.5) {\tiny 2};
  \end{tikzpicture}
}

\def\doubletick{
    \begin{tikzpicture}[scale=0.3]
        \draw[fill=black] (0,.35) -- (.25,0) -- (1,.7) -- (.25, .175) -- cycle;
        \draw[fill=black] (0.5,.35) -- (.75,0) -- (1.5,.7) -- (.75, .175) -- cycle;
        \node[above, white] at (0.35, 0.1) {\tiny 1};
        \node[below, white] at (1.2, 0.5) {\tiny 2};
    \end{tikzpicture}
}

\def\halftick{
    \begin{tikzpicture}[scale=0.3]
        \draw[fill=black] (0,.35) -- (.25,0) -- (1,.7) -- (.25,.15) -- cycle (0.75,0.2) -- (0.77,0.2)  -- (0.6,0.7) -- cycle;
        \node[above] at (0.35, 0.1) {\tiny 1};
        \node[below] at (0.7, 0.5) {\tiny 2};
    \end{tikzpicture}
}

\begin{document}

\begin{center}

\begin{tabularx}{\textwidth}{ XXXX }
  \hline
  \textbf{Subject:}  &  \multicolumn{3}{l}{MODULE NAME} \\
  \textbf{Module Code:} & \multicolumn{3}{l}{MODULE CODE} \\
  \textbf{Time Limit:} & \multicolumn{3}{l}{1 Hour} \\
  \textbf{Marks:} & \multicolumn{3}{l}{\numpoints} \\
  \textbf{Examiner:}  & \multicolumn{3}{l}{Dr SOMEONE} \\
  \hline
\end{tabularx}
\vspace{1cm}
\end{center}

  \noindent\fbox{
    \begin{minipage}[t]{1\columnwidth - 2\fboxsep - 2\fboxrule}
      \textbf{Special Instructions:}
      \begin{itemize}
        \item Answer all questions.
        \item Use of calculator is permitted.
        \item Show your workings.
        \item Round off final answers to four decimal places.
        \item If you choose to write your answers in pencil you then forfeit the right to question the examiners marking.
      \end{itemize}
    \end{minipage}}

  \vspace{1cm}
  \begin{center}
    \pointtable[h][questions]
  \end{center}
  \newpage

\begin{questions}
  \setlength{\rightpointsmargin}{1in}
  \bracketedpoints
  \marksnotpoints
  \totalformat{\fbox{Question \thequestion: \totalpoints}}

  \question[1]What is this?

  \begin{solution}[2.2cm]
  That is a question and this is the solution.
  \end{solution}

  \question

  Suppose that we want to split a question into multiple parts.

  \begin{parts}
  \part[1]
  Is this possible?
  \begin{solution}[1cm]
  Yes!
  \end{solution}

  \part[1]
  Where do the solutions go?
  \begin{solution}[1cm]
  You just place the solution environment below the part.
  \end{solution}

  \end{parts}

\end{questions}

\end{document}
